% Machine-checked proofs of the static model (Lean 4 / Mathlib), added 3 November 2026 at the author's request.
% Source: formal/ (README.md maps every statement to its Lean theorems); built with Lean 4.35.0 and Mathlib v4.35.0.

\section{Machine-checked proofs of the static model}\label{supp:formal}

\emph{What is checked.} The statements of Sections~\ref{sec:model-setup} to~\ref{sec:model-statics} of the main paper are formalized in Lean~4 \citep{moura2021lean4} with the Mathlib library \citep{mathlib2020}, and their proofs are checked by Lean's kernel: Assumptions~\ref{ass:market} to~\ref{ass:capital}, Definition~\ref{def:frontier} and the constants \eqref{eq:constants}; Lemmas~\ref{prop:frontier}, \ref{prop:score}, \ref{prop:saturation} and~\ref{prop:envelope}; \autoref{prop:budget} in discrete and continuous time; Corollaries~\ref{cor:transient} and~\ref{cor:capital}; and parts (i), (ii), (iii) and (v) of \autoref{cor:priorities}. The formalization has \FormalNThm{} main theorems, named after the paper's labels; none contains an unproved step, and each depends only on the three axioms standard in Lean's mathematics (propositional extensionality, choice and quotient soundness). The feasible set, dominance and the frontier are defined as in Definition~\ref{def:frontier}, and the frontier is derived from dominance, not defined as the graph of $m^*$. The code is in the folder \texttt{formal/} of the replication package, whose README maps every part of every statement to its Lean theorems and states each difference of form.

\emph{What is not checked.} Part (iv) of \autoref{cor:priorities}, which rests on \autoref{lem:prices} and \autoref{prop:entry}; the numerical examples inside the statements, which \texttt{code/verify\_theory.py} checks numerically; and every result after Section~\ref{sec:model-statics} (the dynamic optimality of \autoref{prop:dynamic}, the valuation, industry, entry and automation results, the accountability floor, the attainable frontier) and the lemmas and remarks of this supplement. These rest on the written proofs, the symbolic and numerical checks of \texttt{verify\_theory.py} and the referees' re-derivations.

\emph{Readings.} Three informal phrases are given a precise reading. ``To first order'' in Corollaries~\ref{cor:capital} and~\ref{cor:priorities} is read as a statement about the derivative of the maximal score along the direction of the change; the bracket \eqref{eq:kl-bracket} is proved exactly. In continuous time, the induced share path of \autoref{prop:budget}(a) is the solution of the share equation $\dot s=s\,\lambda h(s)\,\sigma^\alpha$ in the sense of Carath\'eodory: an absolutely continuous path that satisfies the equation at almost every date, as it must when spend may jump. The formalization constructs that path from its logistic formula, proves that it has these properties for every measurable spend path with $\int_0^T\sigma^\alpha\,dt<\infty$, and proves the bounds of \autoref{prop:budget}(a) and (c) and \autoref{cor:transient}(b) for every such solution, so their uniqueness, which is not formalized, is not needed. Limits as $s\to\bar s$, $s\to0$ and $g\downarrow G$ are one-sided, and the supremum of \autoref{prop:saturation}(c) is over $s\in(0,\bar s)$. Several hypotheses are not needed: the sign of $X$ is never used, and $k\ge0$ is used only for the dominance results (\autoref{prop:frontier}(a) and (c), \autoref{prop:score}(a)); the formula for the maximal score holds for every real $k$.

\emph{What the check changed.} No statement in scope was found false. Two were imprecise and are corrected in the main paper: the cap condition of \autoref{prop:score} referred to $\sigma^*$ of \eqref{eq:sigstar}, a power of a negative number when $X<k$, and now sets $\sigma^*=0$ when $X\le k$; and the moving intercept of \autoref{prop:budget}(c) must be a measurable path for the integral to be defined.

\emph{How it was produced and audited.} An agent on \texttt{claude-fable-5-1} wrote the formalization. An agent on \texttt{claude-opus-5-5} that had not written it audited it twice against the paper's statements. The first audit found that the continuous-time theorems, as first stated, assumed a share path differentiable at every date and so said nothing about spend paths with jumps; they were restated as above. The second audit confirmed the corrections and wrote instances that apply the main theorems at the paper's illustrative values ($\alpha=0.3$, $\ell=0.6$, $G=0.10$, $k=0.10$, $m_0=0.45$, $X\in\{1,2\}$, $\bar s=0.5$) and to a spend path with a jump, which show that their hypotheses can be met; the instances are part of the package (\texttt{formal/Checks/}). The construction of the induced share path and four further links between the comparative statics and the maximal score were added after the second audit, at its suggestion; the lead agent read their statements but they were not audited a third time. A machine check establishes that the proofs prove the Lean statements; that those statements say what the paper says rests on these audits and on the reader's comparison, which the README is written to make easy.
